% =====================================================================
%  FICHE 11 — THEME 5 — Exercices DIFFICILE — Piles, Nernst, Faraday
% =====================================================================
\documentclass[11pt,a4paper]{article}
\usepackage[utf8]{inputenc}
\usepackage[T1]{fontenc}
\usepackage{lmodern}
\usepackage[french]{babel}
\usepackage{amsmath,amssymb}
\usepackage[version=4]{mhchem}
\usepackage{siunitx}
\sisetup{output-decimal-marker={,},per-mode=symbol}
\usepackage{xcolor}
\usepackage{array}
\usepackage{booktabs}
\usepackage{enumitem}
\usepackage[most]{tcolorbox}
\usepackage{tikz}
\usetikzlibrary{arrows.meta,positioning}
\usepackage[a4paper,margin=2.0cm,headheight=26pt,headsep=10pt]{geometry}
\usepackage{fancyhdr}
\usepackage[hidelinks]{hyperref}

\definecolor{primary}{RGB}{146,95,160}
\definecolor{primarydark}{RGB}{92,52,104}
\definecolor{metalcol}{RGB}{150,150,160}
\definecolor{coppercol}{RGB}{184,115,51}
\definecolor{solcu}{RGB}{170,230,200}
\definecolor{solni}{RGB}{235,225,150}
\definecolor{saltbridge}{RGB}{230,175,90}
\definecolor{wirecol}{RGB}{70,70,75}
\newcommand{\NO}{\ensuremath{\mathrm{N.O.}}}
\newcommand{\Ered}{\ensuremath{E_{\mathrm{r}}^{\circ}}}
\newcommand{\fem}{\ensuremath{\Delta E}}
\newcommand{\cc}[1]{\left[#1\right]}

\newcounter{exo}
\newtcolorbox{exo}[1][]{enhanced,breakable,boxrule=0.9pt,arc=2pt,colback=primary!4,
  colframe=primary,fonttitle=\bfseries,coltitle=white,left=3mm,right=3mm,top=2.2mm,bottom=2.2mm,
  title={Exercice~\refstepcounter{exo}\theexo\ifx\relax#1\relax\relax\else\ \textmd{--- \itshape #1}\fi}}
\newtcolorbox{donnees}{enhanced,boxrule=0.7pt,arc=2pt,colback=black!3,colframe=black!45,
  left=3mm,right=3mm,top=2mm,bottom=2mm,fonttitle=\bfseries\small,coltitle=black!70,title={Données}}

\pagestyle{fancy}\fancyhf{}
\fancyhead[L]{\small\textcolor{primarydark}{\textbf{Fiche 11 · Thème 5} — Piles, Nernst, Faraday}}
\fancyhead[R]{\small\textcolor{primarydark}{\textsc{Difficile}}}
\fancyfoot[C]{\small\thepage}
\renewcommand{\headrulewidth}{0.8pt}
\renewcommand{\headrule}{\hbox to\headwidth{\color{primary}\leaders\hrule height \headrulewidth\hfill}}

\begin{document}
\begin{center}
{\LARGE\bfseries\color{primarydark}Thème 5 — Niveau Difficile}\\[2pt]
{\small\itshape Analyser une pile réelle, appliquer Nernst à une pile complète, et quantifier une électrolyse.}\\[-4pt]
{\color{primary}\rule{\textwidth}{1pt}}
\end{center}

\begin{donnees}
\centering\small
$\Ered(\ce{Cu^{2+}}/\ce{Cu})=+0{,}34$\,V ; $\Ered(\ce{Ni^{2+}}/\ce{Ni})=-0{,}26$\,V ;
$\Ered(\ce{Ag+}/\ce{Ag})=+0{,}80$\,V ; $\Ered(\ce{Zn^{2+}}/\ce{Zn})=-0{,}76$\,V.\\
$F=\SI{96485}{C\per mol}$ ; $M(\ce{Cu})=\SI{63.5}{g\per mol}$ ; $\frac{RT}{F}\ln 10 \simeq 0{,}06$~V.
\end{donnees}

\begin{exo}[analyser une pile cuivre--nickel — d'après annales 2020]
On étudie la pile schématisée ci-dessous (électrodes de cuivre et de nickel, pont salin de \ce{KCl}).

\begin{center}
\begin{tikzpicture}[scale=0.8]
  % béchers
  \fill[solcu] (-3.6,-1.6) rectangle (-2.0,0.0);
  \draw[wirecol,line width=1pt] (-3.6,1.2)--(-3.6,-1.6)--(-2.0,-1.6)--(-2.0,1.2);
  \node[primary,font=\scriptsize] at (-2.8,-1.9){\ce{CuSO4}\ \SI{1}{M}};
  \fill[solni] (2.0,-1.6) rectangle (3.6,0.0);
  \draw[wirecol,line width=1pt] (2.0,1.2)--(2.0,-1.6)--(3.6,-1.6)--(3.6,1.2);
  \node[primary,font=\scriptsize] at (2.8,-1.9){\ce{Ni(NO3)2}\ \SI{1}{M}};
  % électrodes
  \fill[coppercol] (-2.9,1.6) rectangle (-2.7,-1.3);
  \node[wirecol,font=\scriptsize\bfseries] at (-2.8,2.0){Cu};
  \fill[metalcol] (2.7,1.6) rectangle (2.9,-1.3);
  \node[wirecol,font=\scriptsize\bfseries] at (2.8,2.0){Ni};
  % pont salin
  \draw[saltbridge,line width=2.4pt] (-2.8,1.6) .. controls (-2.8,2.7) and (2.8,2.7) .. (2.8,1.6);
  \node[saltbridge!35!black,font=\scriptsize] at (0,2.85){pont salin (\ce{KCl})};
  % voltmètre
  \draw[wirecol,line width=1.2pt] (-2.8,1.6)--(-2.8,3.6)--(-0.8,3.6);
  \draw[wirecol,line width=1.2pt] (2.8,1.6)--(2.8,3.6)--(0.8,3.6);
  \draw[wirecol,line width=1.2pt,fill=white] (-0.8,3.1) rectangle (0.8,4.1);
  \node[primarydark,font=\bfseries] at (0,3.6){V};
\end{tikzpicture}
\end{center}

\begin{enumerate}[label=\textbf{\alph*)},leftmargin=2.2em,topsep=2pt,itemsep=2pt]
  \item À partir des potentiels standard, identifie l'\textbf{anode} et la \textbf{cathode}, avec leur signe.
  \item Écris la demi-réaction qui se produit à chaque électrode.
  \item Calcule la f.é.m. standard $\fem^{\circ}$ de la pile.
  \item Indique le sens de circulation des \textbf{électrons} dans le fil, et le sens de migration des
        ions \ce{Cl^-} et \ce{K+} dans le pont salin.
  \item On propose : \textbf{(A)} l'oxydation $\ce{Cu -> Cu^{2+} + 2 e^-}$ a lieu à l'électrode de cuivre ;
        \textbf{(B)} l'électrode de nickel est la cathode ; \textbf{(C)} les électrons circulent du nickel
        vers le cuivre ; \textbf{(D)} les \ce{Cl^-} migrent vers le compartiment du cuivre et les \ce{K+}
        vers celui du nickel. Laquelle est \textbf{correcte} ?
\end{enumerate}
\end{exo}

\begin{exo}[f.é.m. d'une pile hors conditions standard — d'après livre QCM 26]
On réalise la pile $\ce{Zn(s) | Zn^{2+}(0{,}010\,M) || Ag+(0{,}10\,M) | Ag(s)}$ à \SI{25}{\celsius}.
\begin{enumerate}[label=\textbf{\alph*)},leftmargin=2.2em,topsep=2pt,itemsep=2pt]
  \item Identifie l'anode et la cathode.
  \item Calcule le potentiel de la \textbf{cathode} par la formule de Nernst (couple \ce{Ag+}/\ce{Ag}, $n=1$).
  \item Calcule le potentiel de l'\textbf{anode} par Nernst (couple \ce{Zn^{2+}}/\ce{Zn}, $n=2$).
  \item Déduis-en la f.é.m. $\fem = E_{\text{cathode}} - E_{\text{anode}}$ de la pile.
\end{enumerate}
\end{exo}

\begin{exo}[affinage électrolytique du cuivre — loi de Faraday]
Dans une cuve d'électrolyse servant à purifier le cuivre, la cathode reçoit la réduction
$\ce{Cu^{2+} + 2 e^- -> Cu}$. On fait passer un courant $I=\SI{5.0}{A}$.
\begin{enumerate}[label=\textbf{\alph*)},leftmargin=2.2em,topsep=2pt,itemsep=2pt]
  \item Quelle masse de cuivre se dépose à la cathode en $t=\SI{1.0}{h}$ ?
  \item Combien de temps faut-il, à la même intensité, pour déposer \SI{20}{g} de cuivre ?
        (Exprime en secondes puis en heures.)
\end{enumerate}
\end{exo}

\end{document}
